\chapter{Course organization}
\label{cap:organization}

\section{The three-part path}

\begin{center}\small
\begin{tabular}{lp{4cm}p{7cm}}
\toprule
& Content & Learning objectives \\
\midrule
Part I & Modelling & recognise a link between variables (activation, maximum,
big-M, if and only if\ldots) and prove that the model really imposes it \\
Part II & The problems & apply the links to three families of real problems and to the mixed problems,
from the model to Gurobi code \\
Part III & The course & put what was learned to the test with the additional
modelling questions \\
\bottomrule
\end{tabular}
\end{center}

\section{The format of every exercise (and of the exam)}

Every problem of the course --- and every exam question --- follows the same
four-question scheme:
\begin{enumerate}
\item \textbf{Model.} Write the MILP: variables (with their count), objective,
      constraints with their description, and --- if two linked families of
      variables are involved --- explain the link: what is the implication,
      and how the constraint (or the optimum) imposes it, in both directions.
\item \textbf{Instance.} Write the model for a small numerical instance.
\item \textbf{Heuristic.} Design a constructive algorithm that finds a
      feasible solution of the instance, and run it step by step to obtain
      an upper bound (lower bound if the problem is a maximisation).
\item \textbf{Dual.} Write the dual of the LP relaxation\index{LP relaxation}, for the general
      model and for the instance, and hand-build a feasible dual solution to
      obtain the bound from the other side.
\end{enumerate}
The resulting $\lb \le \zmilp \le \ub$ is the thread running through the
course: a model is not just written down, it is squeezed from both sides
before it is handed to the solver.

\section{The additional questions and their solutions}
Every problem closes with two or three \emph{additional modelling questions}:
they change a datum or add a constraint, and ask how the model changes and what
the new optimum is. They are the part worth practising on, because the base
model is read, the variant is written.

\textbf{The solutions of the additional questions are reserved for
instructors}: they live in a separate booklet, which is not published. What
stays available to everyone --- and it is what one needs in order to know what
is expected --- is the \emph{sandwich worked out in full} on one variant of each
exercise of chapters ``Assignment and scheduling'' and ``Location and covering'': the
model of the variant, a feasible constructive heuristic run step by step, a dual
certificate built by hand and the bound table. It is the model answer the exam
expects.

\section{Grading criteria}

\begin{center}\small
\begin{tabular}{lr}
\toprule
Dimension & Weight \\
\midrule
Correctness of the model (variables, objective, constraints) & 35\% \\
Proof of the link between the variables (both directions) & 25\% \\
Constructive heuristic and correct execution & 20\% \\
Dual of the relaxation and feasible dual solution & 20\% \\
\bottomrule
\end{tabular}
\end{center}

\section{Typical discussion questions}
\begin{itemize}
\item Is the link constraint aggregated or disaggregated? What difference
      does it make to the LP relaxation\index{LP relaxation}?
\item Is the opposite direction of the implication imposed by the constraint
      or does it follow from the optimum? How is it proved?
\item What is the smallest big-M\index{big-M} that can be justified from the data?
\item Does the heuristic find the optimum? How can one know without the
      solver?
\item Is the hand-built dual optimal for the relaxation, or only feasible?
\end{itemize}

\section{The most common mistakes}
\begin{enumerate}
\item Writing an implication whose thesis is true regardless (vacuous
      contrapositive): before writing it, check that antecedent and
      consequent are both genuine.
\item Proving only one direction of an implication ``imposed by the
      constraint'' --- both ``if''s are always needed.
\item Confusing a relation ``imposed by the constraint'' with one that
      ``follows from the optimum'': the second requires the six-step
      exchange argument\index{exchange argument}, not just ``it's clearly worth it''.
\item Concluding ``in every optimal solution'' when the coefficient is only
      $\ge 0$ (not $>0$): the correct conclusion is weaker, ``there exists an
      optimum in which\ldots''.
\item Using a huge big-M\index{big-M} ``to be safe'': it needlessly worsens the LP
      relaxation.
\item Forgetting that a mutual-exclusion constraint ($A+B\le 1$) does not
      imply that at least one of the two must equal 1.
\item Writing the dual of the relaxation with the wrong direction or sign of
      a dual variable relative to the direction of the primal constraint.
\item Building a feasible dual solution without checking its feasibility on
      all the constraints.
\end{enumerate}

\section{Reproducibility}
\begin{lstlisting}[style=output]
python3 -m pip install gurobipy matplotlib pandas
python3 python/run_all.py             # regenerates data, results, figures and notebooks
python3 python/check_numbers.py       # checks every number quoted in the notes
\end{lstlisting}
